\begin{frame}[fragile]\frametitle{Q1(1) 問題}
  \begin{alertblock}{問い}
    任意の$\lambda$-term $M$ と任意の変数$x$について，$M$の帰納法によって，
    $(\lambda x.M)x =_{\beta} M$ を示せ．
  \end{alertblock}
  \begin{itemize}
    \item \textbf{証明のキーポイント:} \alert{帰納法の適用}
    \begin{itemize}
        \item $M$の構造に基づいて帰納法を適用する．
        \item 最小の事象を確認する．
        \item ある事象が成り立つと仮定したとき，一つ大きい事象について成り立つことを示す．
    \end{itemize}
\end{itemize}
\end{frame}

\begin{frame}
\frametitle{Q1 (1) 証明の概要}
\begin{itemize}
    \item 最小の事象について
    \begin{itemize}
        \item $M$ が変数 $x$ のとき
        \item $M$ が定数のとき
    \end{itemize}
    \item ある事象が成り立つと仮定したとき，一つ大きい事象について成り立つことを示す
    \begin{itemize}
        \item $M$ が $M_1 M_2$ のとき
        \item $M$ が $\lambda x. M_1$ のとき
        \item $M$ が $\lambda y. M_1$ ($y \neq x$) のとき
    \end{itemize}
\end{itemize}
\end{frame}

\begin{frame}
\frametitle{Q1 (1) 証明}
\begin{enumerate}
    \item $M \equiv x$のとき
    \begin{align*}
    (\lambda x.x)x &\rhd_{1\beta} [x/x]x \\
    &\equiv x \\
    &\equiv M \\
    \end{align*}

    \item $M \equiv a$かつ$a \not\equiv x$のとき
    \begin{align*}
    (\lambda x.a)x &\rhd_{1\beta} [x/x]a \\
    &\equiv a \\
    &\equiv M \\
    \end{align*}
\end{enumerate}
\end{frame}

\begin{frame}
    \frametitle{Q1 (1) 証明}
    \begin{enumerate}
        \item[3] $M \equiv M_1 M_2$のとき
        
        $(\lambda x.M_1)x =_{\beta} M_1$ と $(\lambda x.M_2)x =_{\beta} M_2$ 
        が成り立つと仮定する．
        \begin{align*}
        (\lambda x.M)x &\equiv (\lambda x.(M_1 M_2))x \\
        &\rhd_{1\beta} [x/x](M_1 M_2) \\
        &\equiv ([x/x]M_1 [x/x]M_2) \\
        (\lambda x.M_1)x \rhd_{1\beta} [x/x]M_1 \text{より}=_{\beta}\text{の定義から}\\
        &=_{\beta}((\lambda x.M_1)x ((\lambda x.M_2)x)) \\
        &=_{\beta} (M_1 ((\lambda x.M_2)x)) &\text{（仮定より）}\\
        &=_{\beta} M_1 M_2  &\text{（仮定より）}\\
        &\equiv M \\
        \end{align*}
    \end{enumerate}
\end{frame}

\begin{frame}
    \frametitle{Q1 (1) 証明}
    \begin{enumerate}
        \item[4] $M \equiv \lambda x.M_1$のとき
        \begin{align*}
            (\lambda x.M)x &\equiv (\lambda x.(\lambda x.P))x\\
            &\rhd_{1\beta} [x/x](\lambda x.P) \\
            &\equiv \lambda x.P \\
            &\equiv M \\
        \end{align*}
    \end{enumerate}
\end{frame}

\begin{frame}
    \frametitle{Q1 (1) 証明}
    \begin{enumerate}
        \item[5] $M \equiv \lambda y.M_1$ ($y \neq x$) のとき
        \begin{align*}
            (\lambda x.M)x &\equiv (\lambda x.(\lambda y.M_1))x \\
            &\rhd_{1\beta} [x/x](\lambda y.M_1) \\
            &\equiv \lambda y.[x/x]M_1 \\
            &=_{\beta} (\lambda y.(\lambda x.M_1)x) \\
            &=_{\beta} \lambda y.(M_1) &\quad\text{（仮定より）}\\
            &\equiv M \\
        \end{align*}
    \end{enumerate}
    上記の1$\sim$5より，$(\lambda x.M)x =_{\beta} M$が示された． $\square$
\end{frame}

\begin{frame}
    \frametitle{Q1 (2) 問題}
    \begin{alertblock}{問い}
    任意の$\lambda$-term $Z$ および変数$x,y_1,y_2,y_3$について，
    \begin{align*}
        Y(\lambda xy_1y_2y_3.Z)y_1y_2y_3 =_{\beta} [(Y(\lambda xy_1y_2y_3.Z))/x]Z \\
    \end{align*}
    であることを証明せよ． ここで$Y \equiv (\lambda ux.x(uux))(\lambda ux.x(uux))$とする．

    term $Y$は固定点コンビネータであり、すべてのterm $X$に対して$YX =_{\beta} X(YX)$が成り立つ．
    \end{alertblock}
    \begin{itemize}
    \item \textbf{証明のポイント:} 
        \begin{itemize}
        \item 左辺を展開して右辺と同様になることを示す．
        \end{itemize}
    \end{itemize}
\end{frame}

\begin{frame}
    \frametitle{Q1 (2) 証明}
    左辺の省略形を直す．
    \begin{align*}
        Y(\lambda xy_1y_2y_3.Z)y_1y_2y_3 &\equiv ((((Y(\lambda x.(\lambda y_1.(\lambda y_2.(\lambda y_3.Z)))))y_1)y_2)y_3) \\
    \end{align*}
    $YX =_{\beta} X(YX)$ より
    \begin{align*}
        &=_{\beta} (((((\lambda x.(\lambda y_1.(\lambda y_2.(\lambda y_3.Z))))(Y(\lambda x.(\lambda y_1.(\lambda y_2.(\lambda y_3.Z))))))y_1)y_2)y_3) \\
    \end{align*}
            ここで $Y(\lambda x.(\lambda y_1.(\lambda y_2.(\lambda y_3.Z))))$ を $W$ とする．
    \begin{align*}
        &\equiv (((((\lambda x.(\lambda y_1.(\lambda y_2.(\lambda y_3.Z))))(W))y_1)y_2)y_3) \\
    \end{align*}
\end{frame}

\begin{frame}
    \frametitle{Q1 (2) 証明}
    \begin{align*}
        &\rhd_{1\beta} (((([W/x](\lambda y_1.(\lambda y_2.(\lambda y_3.Z))))y_1)y_2)y_3) \\
        &\equiv ((((\lambda y_1.(\lambda y_2.(\lambda y_3.([W/x]Z))))y_1)y_2)y_3) \\
        &\rhd_{1\beta} (([y_1/y_1](\lambda y_2.(\lambda y_3.([W/x]Z)))y_2)y_3) \\
        &\equiv (((\lambda y_2.(\lambda y_3.([W/x]Z)))y_2)y_3) \\
        &\rhd_{1\beta} ([y_2/y_2](\lambda y_3.([W/x]Z))y_3) \\
        &\equiv ((\lambda y_3.([W/x]Z))y_3) \\
        &\rhd_{1\beta} ([y_3/y_3]([W/x]Z)) \\
        &\equiv ([W/x]Z) \\
        W \text{を直すと}\\
        &\equiv [(Y(\lambda xy_1y_2y_3.Z))/x]Z \qquad \square\\
    \end{align*}
    
\end{frame}